\begin{frame}{State를 먼저 정의하라}
\dpstate{\texttt{dp[i][j]}가 무엇을 뜻하는지 완전한 문장으로 말할 수 있어야 한다.}
\begin{alertblock}{순서}DP table을 먼저 만들지 말고 state의 의미를 먼저 정의한다.\end{alertblock}
\end{frame}
\begin{frame}{DP 핵심 설계 Template}
\small
\dpdesign
{부분 문제의 의미를 완전한 문장으로 정의}
{더 작은 dependency state로 recurrence 구성}
{더 이상 transition을 적용하지 않는 경계 정의}
{모든 dependency가 먼저 계산되는 순서 선택}
{원문제의 값이 저장되는 state 명시}
\end{frame}
\begin{frame}{Template 다음의 Engineering Check}
\begin{enumerate}
\item dependency graph와 evaluation order가 일치하는가?
\item 값뿐 아니라 실제 solution reconstruction이 필요한가?
\item computed state 수와 state당 transition 비용은 얼마인가?
\item 기본 table과 recursion stack의 공간은 얼마인가?
\item space optimization 후에도 answer와 복원이 가능한가?
\end{enumerate}
\end{frame}
\begin{frame}{복잡도는 State 수 × Transition 비용}
\[
\text{time}=(\text{계산되는 state 수})\times(\text{state당 transition 비용})
\]
\begin{exampleblock}{예}Matrix Path는 $nm$개 state, 각 state에서 두 predecessor 비교: $\Theta(nm)$.\end{exampleblock}
\end{frame}
\begin{frame}{설계 오류 Checkpoint}
\begin{itemize}\item 의미 없는 \texttt{dp[i][j]}\item base case 누락\item dependency보다 먼저 계산\item answer cell을 모름\item length만 구하고 실제 해를 복원하지 못함\end{itemize}
\dptakeaway{표의 모양보다 state의 의미와 dependency graph가 먼저다.}
\vfill\muted{이제 DP가 정답을 보존하는 조건을 정리한 뒤 같은 template를 실제 문제에 적용한다.}
\end{frame}
